\documentclass[11pt]{article}
\usepackage{mathblatt}
% Präambel des Setzers (werkzeuge/setzer-praeambel.tex, wortgleich); dazu \gteil (Muster T2U)
\makeatletter
\setlength{\mbpfbe}{0mm}% keine Punkte (pruefung.md § 5: nur beim Prüfstein)
\newlength{\gEin}\setlength{\gEin}{0pt}
\let\gSchrift\relax
\renewcommand{\pfkreuzzeile}[1]{\par\noindent\hangindent3.2em\hangafter1\hspace*{1.6em}$\square$\ #1\par}
\newcommand{\gkreuz}[1]{$\square$\ #1\hspace{2em}}
\newcommand{\gkreuzzeile}[2]{\par\noindent\hangindent2.2em\hangafter1\mbox{#1\,$\square$}\ #2\par}% Buchstabe, Option: eine je Zeile (neben Skizze)
% links, was man ansieht, rechts, was man tut; Spaltengrenze überall an derselben Stelle (gemeinsam.md „Satz“)
\newsavebox{\gbox}\newlength{\gLinks}
% v0.10 (Lehrer 13.10.): keine starre Halbierung – die linke Spalte ist so breit wie die Figur (+ 6 mm), die
% Antwort steht direkt daneben; ist die Figur breiter als 0,62 Zeile, steht die Antwort darunter
\newcommand{\gzwei}[2]{\par\smallskip\noindent\sbox{\gbox}{#1}%
  \setlength{\gLinks}{\dimexpr\wd\gbox+6mm\relax}%
  \ifdim\gLinks>0.62\linewidth
    \usebox{\gbox}\par\smallskip\noindent\begin{minipage}{\linewidth}#2\end{minipage}%
  \else
    \begin{minipage}[t]{\gLinks}\vspace{0pt}\usebox{\gbox}\end{minipage}%
    \begin{minipage}[t]{\dimexpr\linewidth-\gLinks\relax}\vspace{0pt}\raggedright #2\end{minipage}%
  \fi\par}
% v0.9 (12.10., pruefung.md § 5 „Marke mit B, Bedingungen rechts“): drittes Argument von pfaufg = Bedingungen
% der Prüfung (v0.10: \pfohnetr = TR durchgestrichen, „FOR“), klein grau im rechten Rand auf der Grundlinie der Nummer, Zeilen mit \\
\newcommand{\gbedrand}[2]{\ifblank{#2}{}{\rlap{\hspace*{\dimexpr#1+2mm\relax}%
  \parbox[t]{15mm}{\raggedright\scriptsize\leavevmode\color{mbgrau}\strut#2}}}}
% Teil eines Bündels (a, b …): Marke links im Rand, Bedingungen rechts im Rand, auf der Zeile des Teils
\newcommand{\gteilrand}[2]{\llap{\mbpfrandmarke{#1}\hspace*{\mbpfnr}}\gbedrand{\linewidth}{#2}}
% Teil eines Bündels: Buchstabe in eigener Spalte, Inhalt eingerückt (\gEin = Einzug, damit die Spaltengrenze von
% \gzwei an derselben Stelle bleibt)
\newenvironment{gbteil}[3]{\par\smallskip\noindent\gteilrand{#2}{#3}\makebox[\mbpfteil][l]{\textbf{#1}}%
  \begin{minipage}[t]{\dimexpr\linewidth-\mbpfteil\relax}\setlength{\gEin}{\mbpfteil}\raggedright\strut\ignorespaces}%
 {\end{minipage}\par}
% Grafik höchstens so breit wie die Spalte (Bündel nebeneinander, Koordinatensystem ohne Karo)
\newsavebox{\gpbox}
\newcommand{\gpassend}[1]{\sbox{\gpbox}{#1}\ifdim\wd\gpbox>\linewidth\typeout{GPASSEND: verkleinert \the\wd\gpbox\space auf \the\linewidth}\resizebox{\linewidth}{!}{\usebox{\gpbox}}\else\usebox{\gpbox}\fi}
% Skizzenfeld (v0.10, Lehrer 13.10.): Karo ab dem linken Rand über die Zeile, „Skizze“ klein grau oben links im Feld
% (wie \pfskizzenfeld der Vorlage); der rechte Rand bleibt den Bedingungen
\newcommand{\gskizze}[1]{\par\smallskip\noindent\begin{tikzpicture}
  \clip (0,0) rectangle ({\linewidth-1mm},{#1*10mm});
  \draw[black!16,line width=0.3pt,step=5mm] (0,0) grid ({\linewidth},{#1*10mm});
  \node[anchor=north west,inner sep=2pt,font=\scriptsize,text=mbgrau,fill=white] at (0,{#1*10mm}) {Skizze};
  \end{tikzpicture}\par}
\RenewDocumentEnvironment{pfaufg}{m m m}{%
  \begin{lrbox}{\mbaufgabebox}\begin{minipage}[t]{\linewidth}%
  \gSchrift\noindent\hspace*{\gEin}\mbpfrandmarke{#1}\gbedrand{\linewidth-\gEin}{#3}%
  \begin{minipage}[t]{\mbpfnr}\strut\textbf{#2}\end{minipage}%
  \begin{minipage}[t]{\dimexpr\linewidth-\gEin-\mbpfnr\relax}\raggedright\strut\ignorespaces}%
 {\end{minipage}%
  \end{minipage}\end{lrbox}%
  \setlength{\mbaufgabehoehe}{\dimexpr\ht\mbaufgabebox+\dp\mbaufgabebox\relax}%
  \ifdefempty{\gkopfnach}{\par\addvspace{5pt}\Needspace*{\mbaufgabehoehe}}{%
    % Stufenkopf, der mit dieser (ersten) Aufgabe zusammenbleibt: Platz für beide verlangen, dann den Kopf setzen
    % (12.10., Quadratische: ein Kopf stand allein unten, die Aufgabe mit Koordinatensystem auf der nächsten Seite)
    \par\Needspace*{\dimexpr\mbaufgabehoehe+40pt\relax}\gkopfnach\gdef\gkopfnach{}\par\addvspace{5pt}}%
  \noindent\usebox{\mbaufgabebox}\par\addvspace{5pt}}
% Stufenkopf vormerken: \gkopfvor{<Kopf>} – die nächste pfaufg setzt ihn, sobald beide zusammen Platz haben
\def\gkopfnach{}
\newcommand{\gkopfvor}[1]{\gdef\gkopfnach{#1}}
% Teilaufgabe mit Skizze links in gzwei: Buchstabe in eigener Spalte, darunter Text und Skizze
\newcommand{\gteil}[2]{\makebox[\mbpfteil][l]{\textbf{#1}}\parbox[t]{\dimexpr0.5\textwidth-\gEin-\mbpfnr-5mm-\mbpfteil\relax}{\raggedright #2}}
% Stufenkopf der Vorlage (mathblatt.sty 2026-10-12: Linie, Abstand, 13 pt, \pfjahre) ohne ihren Vorhalt von
% 14 Zeilen: der Setzer fordert vor dem Kopf 0,16 Seitenhöhe (Kopf + erste Aufgabe), mehr nicht
\setlength{\mbpfstufevorhalt}{0pt}
\makeatother
\newcommand{\karo}[1]{\par\smallskip\noindent\begin{tikzpicture}
  \clip (0,0) rectangle ({\linewidth-1mm},{#1*10mm});
  \draw[black!16,line width=0.3pt,step=5mm] (0,0) grid ({\linewidth},{#1*10mm});
  \end{tikzpicture}\par}
\newcommand{\kk}[1]{$\square$\,#1\hspace{0.9em}}
\newcommand{\linie}[1]{\,{\color{black!45}\rule[-1pt]{#1}{0.4pt}}\,}
\newcommand{\pkt}{\textperiodcentered{}}
% v0.14 (14.10.): Stufenkopf mit fünf Punkten (Option = Jahre für die Bestellkarte) oder „zuletzt 19“
\NewDocumentCommand{\gpunkte}{o m}{#2}
\newcommand{\gzuletzt}[1]{zuletzt #1}
% v0.7 (12.10.): Schülerfrage unter dem Stufenkopf; Ankreuzoption als Bild (Buchstabe und Kästchen darunter)
\newcommand{\gfrage}[1]{\par\nobreak\noindent\hspace*{\mbpfnr}{\small\itshape #1}\par\nobreak\addvspace{2pt}}
\newlength{\gbildbreite}\setlength{\gbildbreite}{0.32\linewidth}% Setzer setzt sie je Reihe (vier Bilder: 0,24)
\newcommand{\gbildoption}[3]{\begin{minipage}[t][#3][b]{\gbildbreite}\centering #2\par\smallskip\mbox{#1\,$\square$}\end{minipage}}
% Lineare (12.10.): Bild mit Buchstabe ohne Kästchen (Diagramm zuordnen); Aussage mit „wahr / falsch“ zum Ankreuzen
\newcommand{\gbildname}[3]{\begin{minipage}[t][#3][b]{\gbildbreite}\centering #2\par\smallskip\mbox{\textbf{#1}}\end{minipage}}
\newcommand{\gwahrfalsch}[2]{\par\noindent\hangindent2.2em\hangafter1\makebox[2.2em][l]{(#1)}#2\hfill\mbox{wahr $\square$\quad falsch $\square$}\par}
\begin{document}

\pfheftstil
\pflaufkopf[Abitur GK \pkt{} Prüfung]{Lage und Schnittpunkt}{AG50}\pfseitevon
\pfheftkopf{Lage und Schnittpunkt}{}
\pfstufenzeile{Lage begründen \pkt{} Schnittpunkt gesucht}

% ===================== Stufe: Lage begründen =====================
\gkopfvor{\pfstufekopf{Lage begründen}{\gpunkte[24 23 22]{\textcolor{black!25}{$\circ$}\,\textcolor{black!25}{$\circ$}\,\textcolor{black!80}{$\bullet$}\,\textcolor{black!80}{$\bullet$}\,\textcolor{black!80}{$\bullet$}}}\special{pdf:dest (stufe-lage-begrunden) [@thispage /XYZ @xpos @ypos null]}}
% kein Einstieg: keine Schülerfrage in der Gliederung (Schülerfrage, -Fälle, -Form, -Weg, -Skizze; oder Frage) – Einstieg nur von dort, nie aus der Bank (pruefung.md § 5, 12.10.)
% 2022-bebb-gk-A1.4b-m1 (Original 2022-bebb-gk-A1.4b, Schwere I, rang 1; herausgelöste Fassung: keine)
\begin{pfaufg}{GK~’22\\A1.4b}{1.}{\pfohnetr}
Gegeben ist die Ebene $E\colon 2x - 5y = 0$. Beschreibe die Lage von $E$ im Koordinatensystem.
\hfill E enthält die z-Achse $=$\linie{30mm}
\fusshilfe{\mbox{1:~E enthält die z-Achse}}
\end{pfaufg}
% 2024-bebb-gk-A1.2a-m1 (Original 2024-bebb-gk-A1.2a, Schwere I, rang 2; herausgelöste Fassung: keine)
\begin{pfaufg}{GK~’24\\A1.2a}{2.}{\pfohnetr}
Gegeben sind $P(1 \mid 0 \mid 17)$ und $Q(5 \mid t \mid 13)$ mit $t \in \mathbb{R}$. Entscheide, ob es ein $t$ gibt, für das die Gerade PQ parallel zur xy-Ebene verläuft. Begründe.
\karo{2}
\fusshilfe{\mbox{2:~nein}}
\end{pfaufg}
% 2022-bebb-gk-B3e-m1 (Original 2022-bebb-gk-B3e, Schwere I, rang 3; herausgelöste Fassung: keine)
\begin{pfaufg}{GK~’22\\B3e}{3.}{}
Ein Kirchturmdach steht auf dem Quadrat ABCD mit $A(0 \mid 0 \mid 0)$, $B(6 \mid 0 \mid 0)$, $C(6 \mid 6 \mid 0)$, $D(0 \mid 6 \mid 0)$. Darüber liegen $E(3 \mid 0 \mid 5)$, $F(6 \mid 3 \mid 5)$, $G(3 \mid 6 \mid 5)$, $H(0 \mid 3 \mid 5)$ und die Spitze $S(3 \mid 3 \mid 10)$; 1 LE entspricht 1 m. Das Dach hat vier gleiche viereckige Dachflächen wie CGSF und vier dreieckige Giebelflächen. Gegeben sind die Ebenen $M_1\colon x = 6$, $M_2\colon x - y = 0$ und $M_3\colon z = 5$. Eine davon ist eine Symmetrieebene des Dachs. Gib sie an und beschreibe ihre Lage.
\gzwei{\begin{ksys3}[xyz,x1min=0,x1max=6,x2min=0,x2max=6,x3min=0,x3max=10,karo=0.16]\rvieleck{6,0,0;6,6,0;6,3,5}\rvieleck{6,6,0;0,6,0;3,6,5}\rstrecke[dashed]{0,0,0}{6,0,0}\rstrecke{6,0,0}{6,6,0}\rstrecke{6,6,0}{0,6,0}\rstrecke[dashed]{0,6,0}{0,0,0}\rstrecke[dashed]{0,0,0}{3,0,5}\rstrecke{3,0,5}{6,0,0}\rstrecke{6,0,0}{6,3,5}\rstrecke{6,3,5}{6,6,0}\rstrecke{6,6,0}{3,6,5}\rstrecke{3,6,5}{0,6,0}\rstrecke{0,6,0}{0,3,5}\rstrecke[dashed]{0,3,5}{0,0,0}\rstrecke{3,0,5}{3,3,10}\rstrecke{6,3,5}{3,3,10}\rstrecke{3,6,5}{3,3,10}\rstrecke{0,3,5}{3,3,10}\rpunkt*[above left]{0}{0}{0}{A}\rpunkt*[below left]{6}{0}{0}{B}\rpunkt*[below right]{6}{6}{0}{C}\rpunkt*[above right]{0}{6}{0}{D}\rpunkt*[left]{3}{0}{5}{E}\rpunkt*[below right]{6}{3}{5}{F}\rpunkt*[right]{3}{6}{5}{G}\rpunkt*[above left]{0}{3}{5}{H}\rpunkt*[above]{3}{3}{10}{S}\end{ksys3}}{\karo{2}}
\fusshilfe{\mbox{3:~sie enthält die z-Achse und die Diagonale AC und steht senkrecht auf der Grundfläche}}
\end{pfaufg}
% 2023-bebb-gk-B3b-m1 (Original 2023-bebb-gk-B3b, Schwere II, rang 4; herausgelöste Fassung: keine)
\begin{pfaufg}{GK~’23\\B3b}{4.}{}
Die Punkte $A(6 \mid 0 \mid 0)$, $B(0 \mid 6 \mid 0)$ und $C(0 \mid 0 \mid 6)$ liegen in der Ebene $L_1\colon x + y + z = 6$, die Punkte $D(3{,}5 \mid 0 \mid 0)$, $E(0 \mid 3{,}5 \mid 0)$ und $F(0 \mid 0 \mid 3{,}5)$ in der Ebene $L_2\colon 2x + 2y + 2z = 7$. Sie sind die Ecken des Körpers ABCDEF. Begründe, dass die Ebenen $L_1$ und $L_2$ keinen gemeinsamen Punkt haben.
\gzwei{\begin{ksys3}[xyz,x1min=0,x1max=6,x2min=0,x2max=6,x3min=0,x3max=6,karo=0.2]\rvieleck{6,0,0;0,6,0;0,0,6}\rstrecke{6,0,0}{3.5,0,0}\rstrecke{0,6,0}{0,3.5,0}\rstrecke{0,0,6}{0,0,3.5}\rstrecke[dashed]{3.5,0,0}{0,3.5,0}\rstrecke[dashed]{0,3.5,0}{0,0,3.5}\rstrecke[dashed]{0,0,3.5}{3.5,0,0}\rpunkt*[left]{6}{0}{0}{A}\rpunkt*[below]{0}{6}{0}{B}\rpunkt*[right]{0}{0}{6}{C}\rpunkt*[below left]{3.5}{0}{0}{D}\rpunkt*[below]{0}{3.5}{0}{E}\rpunkt*[left]{0}{0}{3.5}{F}\end{ksys3}}{\karo{2}}
\fusshilfe{\mbox{4:~L$_1$ $\parallel$ L$_2$; kein gemeinsamer Punkt}}
\end{pfaufg}

% ===================== Stufe: Schnittpunkt gesucht =====================
\gkopfvor{\pfstufekopf{Schnittpunkt gesucht}{\gpunkte[23]{\textcolor{black!25}{$\circ$}\,\textcolor{black!25}{$\circ$}\,\textcolor{black!25}{$\circ$}\,\textcolor{black!80}{$\bullet$}\,\textcolor{black!25}{$\circ$}}}\special{pdf:dest (stufe-schnittpunkt-gesucht) [@thispage /XYZ @xpos @ypos null]}}
% kein Einstieg: keine Schülerfrage in der Gliederung (Schülerfrage, -Fälle, -Form, -Weg, -Skizze; oder Frage) – Einstieg nur von dort, nie aus der Bank (pruefung.md § 5, 12.10.)
% 2023-bebb-gk-B3d-m1 (Original 2023-bebb-gk-B3d, Schwere II, rang 5; herausgelöste Fassung: keine)
\begin{pfaufg}{GK~’23\\B3d}{5.}{}
Die Punkte $A(6 \mid 0 \mid 0)$, $B(0 \mid 6 \mid 0)$ und $C(0 \mid 0 \mid 6)$ liegen in der Ebene $L_1\colon x + y + z = 6$, die Punkte $D(3{,}5 \mid 0 \mid 0)$, $E(0 \mid 3{,}5 \mid 0)$ und $F(0 \mid 0 \mid 3{,}5)$ in der Ebene $L_2\colon 2x + 2y + 2z = 7$. Sie sind die Ecken des Körpers ABCDEF. M(3 | 3 | 0) ist der Mittelpunkt von AB. Die Gerade durch $P(0 \mid 0 \mid 3)$ und $M$ schneidet $L_2$ im Punkt $Q$. Ermittle die Koordinaten von $Q$.
\gzwei{\begin{ksys3}[xyz,x1min=0,x1max=6,x2min=0,x2max=6,x3min=0,x3max=6,karo=0.2]\rvieleck{6,0,0;0,6,0;0,0,6}\rstrecke{6,0,0}{3.5,0,0}\rstrecke{0,6,0}{0,3.5,0}\rstrecke{0,0,6}{0,0,3.5}\rstrecke[dashed]{3.5,0,0}{0,3.5,0}\rstrecke[dashed]{0,3.5,0}{0,0,3.5}\rstrecke[dashed]{0,0,3.5}{3.5,0,0}\rpunkt*[left]{6}{0}{0}{A}\rpunkt*[below]{0}{6}{0}{B}\rpunkt*[right]{0}{0}{6}{C}\rpunkt*[below left]{3.5}{0}{0}{D}\rpunkt*[below]{0}{3.5}{0}{E}\rpunkt*[left]{0}{0}{3.5}{F}\end{ksys3}}{\karo{2}}
\fusshilfe{\mbox{5:~$\left(\begin{smallmatrix}0\\0\\3\end{smallmatrix}\right)$ + t · $\left(\begin{smallmatrix}3\\3\\-3\end{smallmatrix}\right)$; Q(0,5|0,5|2,5)}}
\end{pfaufg}
% 2023-bebb-gk-A1.4b-m1 (Original 2023-bebb-gk-A1.4b, Schwere III, rang 6; herausgelöste Fassung: keine)
\begin{pfaufg}{GK~’23\\A1.4b}{6.}{\pfohnetr}
Gegeben ist die Gerade $g\colon \vec{x} = \begin{pmatrix}1\\4\\-3\end{pmatrix} + s \cdot \begin{pmatrix}2\\0\\3\end{pmatrix}$, $s \in \mathbb{R}$. Die Gerade $h$ verläuft durch $A(5 \mid 0 \mid 0)$ und $B(6 \mid 1 \mid b)$ mit $b \in \mathbb{R}$. $g$ und $h$ haben einen gemeinsamen Punkt. Ermittle $b$.
\karo{2}
\fusshilfe{\mbox{6:~$\left(\begin{smallmatrix}5\\0\\0\end{smallmatrix}\right)$ + r · $\left(\begin{smallmatrix}1\\1\\b\end{smallmatrix}\right)$; 4; Schnittpunkt (9|4|9)}}
\end{pfaufg}

\pfblattende{10}{26 & 25 & 24 & 23 & 22 & 21 & 20 & 19 & 18 & 17}{ &  & A1.2a & A1.4b B3b B3d & A1.4b B3e & A1.4b & B3.1d & B3.1c &  & }%
{\pfausschnitt{Geraden und Ebenen}{Lage und Schnittpunkt}{}}
\end{document}